\chapter{Modularization and ADTs}

\textbf{Readings:} CP:AMA 15 (writing large programs), 19.1--19.3.

\section{Modularization}

\begin{definition}
\emph{Modularization} is dividing a program into smaller, well-defined pieces called
\emph{modules}. A module provides a collection of functions (and possibly other things, such as
constants or data structures) that share a common aspect or purpose.
\end{definition}

We have been doing a small version of this since CS 135 with helper functions. A module is the
next level up: a whole group of related functions.

\paragraph{Modules vs.\ files.} It is good style to store each module in its own file
(in C, two files: \code{.h} and \code{.c}). In these notes ``module'' and ``file'' are used
almost interchangeably.

\subsection{Terminology}

\begin{itemize}
  \item A \emph{client} is a module (or program) that requires a function a module provides.
  \item A module may itself be a client of other modules. The ``requires'' relationships form a
    \emph{module dependency graph}.
  \item The dependency graph must \textbf{not have cycles} (module A cannot require B if B
    requires A).
  \item There must be a \emph{root} (main) file that is only a client and is not required by
    anything; in C it is the file that contains \code{main}.
\end{itemize}

\begin{center}
\begin{tikzpicture}[m/.style={draw,rounded corners,fill=blue!6,font=\small\ttfamily,minimum height=7mm,minimum width=2.2cm},
  arr/.style={-{Stealth},thick}]
  \node[m,fill=orange!15] (main) at (0,0) {main.c};
  \node[m] (dict) at (-2.6,-1.4) {dictionary};
  \node[m] (io) at (2.6,-1.4) {input};
  \node[m] (bst) at (-2.6,-2.8) {bst};
  \node[m] (str) at (2.6,-2.8) {strings};
  \draw[arr] (main) -- (dict);
  \draw[arr] (main) -- (io);
  \draw[arr] (dict) -- (bst);
  \draw[arr] (io) -- (str);
  \draw[arr] (dict) -- (str);
  \node[font=\scriptsize,anchor=west,align=left] at (4.1,-0.1) {root: only a client};
  \node[font=\scriptsize,anchor=west,align=left] at (4.1,-2.15) {arrow: ``requires''\\(no cycles allowed)};
\end{tikzpicture}
\end{center}

\subsection{Motivation}

\begin{enumerate}
  \item \textbf{Re-usability.} A well-written module (e.g.\ a dictionary) can be used by many
    programs.
  \item \textbf{Maintainability.} Bugs can be found and fixed in one place; modules can be
    tested independently; big projects can be split among team members.
  \item \textbf{Abstraction.} A client only needs to know \emph{what} a module does, not
    \emph{how} it does it.
\end{enumerate}

In computer science, modularization is also known as the \emph{separation of concerns} (SoC).

\section{Modules in Racket}

\begin{lstlisting}[style=racket]
;; fun.rkt
(provide fun?)

(define lofn '(-3 7 42 136 1337 4010 8675309))

;; (fun? n) determines if n is a fun integer
;; fun?: Int -> Bool
(define (fun? n)
  (not (false? (member n lofn))))
\end{lstlisting}
The \code{provide} special form makes \code{fun?} available to clients; \code{lofn} is not
provided, so it is private to the module.
\begin{lstlisting}[style=racket]
;; client.rkt
(require "fun.rkt")
(fun? 4010)    ; => #t
(fun? 4011)    ; => #f
\end{lstlisting}

\section{Scope Revisited}

\begin{itemize}
  \item \textbf{Local scope:} the identifier is only visible inside the local region (e.g.\ the
    function body or block) where it is defined.
  \item \textbf{Global scope:} the identifier is defined at the top level and is visible to all
    code that follows its definition.
\end{itemize}
Modules split global scope into two levels:
\begin{itemize}
  \item \textbf{Module scope:} visible only inside the module it is defined in.
  \item \textbf{Program scope:} visible outside the module (to clients) as well.
\end{itemize}
In Racket, identifiers are module scope unless \code{provide}d. In C it is the other way round.

\begin{important}[Scope in C]
Global identifiers in C have \emph{program} scope by default. Adding the keyword
\code{static} to a global definition gives it \emph{module} scope:
\begin{lstlisting}[style=c]
static const int secret = 42;     // module scope
static int helper(int n) { ... }  // module scope ("private" helper)
int visible(int n) { ... }        // program scope
\end{lstlisting}
(The keyword \code{static} means something different on local variables; we avoid that use.)
\end{important}

To use a program-scope identifier from another file, the client must \emph{declare} it.
For variables this needs the keyword \code{extern} (\code{extern int counter;}); for functions
a prototype is enough. Of course clients should get these declarations by including the
module's header file rather than writing them by hand.

\section{Interfaces}

\begin{definition}
A module's \emph{interface} is the list of functions (and other identifiers) it provides,
together with their documentation. The \emph{implementation} is the code that defines them.
\end{definition}

Interface documentation should contain:
\begin{itemize}
  \item an overall description of the module;
  \item the list of functions it provides;
  \item the contract and purpose (and requires/effects/time) of each provided function.
\end{itemize}
The interface should \emph{not} describe how things are implemented; that is the point.

\subsection{Example: a sum module in Racket}
\begin{lstlisting}[style=racket]
;; sum.rkt: a module for summing numbers
(provide sum-first sum-squares)

;; (sum-first n) produces 1 + 2 + ... + n
;; sum-first: Nat -> Nat
;; (sum-squares n) produces 1^2 + 2^2 + ... + n^2
;; sum-squares: Nat -> Nat

;;;;;;;;;;;;;;;; IMPLEMENTATION ;;;;;;;;;;;;;;;;
(define (sum-first n)
  (/ (* n (+ n 1)) 2))

(define (sum-squares n)
  (/ (* n (+ n 1) (+ (* 2 n) 1)) 6))

(define (private-helper p) ...)   ; not provided
\end{lstlisting}

\subsection{The same module in C}

\begin{lstlisting}[style=c]
// sum.h -- INTERFACE
// a module for summing numbers

// sum_first(n) produces 1 + 2 + ... + n
// requires: n >= 0
int sum_first(int n);

// sum_squares(n) produces 1^2 + 2^2 + ... + n^2
// requires: n >= 0
int sum_squares(int n);
\end{lstlisting}
\begin{lstlisting}[style=c]
// sum.c -- IMPLEMENTATION
#include "sum.h"

int sum_first(int n) {
  return n * (n + 1) / 2;
}

int sum_squares(int n) {
  return n * (n + 1) * (2 * n + 1) / 6;
}
\end{lstlisting}
\begin{lstlisting}[style=c]
// main.c -- CLIENT (and test client)
#include <assert.h>
#include "sum.h"

int main(void) {
  assert(sum_first(1) == 1);
  assert(sum_first(10) == 55);
  assert(sum_squares(3) == 14);
}
\end{lstlisting}
The implementation includes its own header so the compiler can check that the declarations
and definitions match.

\begin{center}
\begin{tikzpicture}[m/.style={draw,rounded corners,font=\small\ttfamily,minimum height=8mm,minimum width=2.6cm,align=center},
  arr/.style={-{Stealth},thick}]
  \node[m,fill=blue!8] (h) at (0,0) {sum.h\\{\scriptsize\rmfamily interface}};
  \node[m,fill=green!8] (c) at (-3.8,-1.8) {sum.c\\{\scriptsize\rmfamily implementation}};
  \node[m,fill=orange!15] (mn) at (3.8,-1.8) {main.c\\{\scriptsize\rmfamily client}};
  \draw[arr] (c) -- node[above left,font=\scriptsize\ttfamily]{\#include} (h);
  \draw[arr] (mn) -- node[above right,font=\scriptsize\ttfamily]{\#include} (h);
  \draw[dashed,red!60!black,thick] (0,-0.9) -- (0,-2.6);
  \node[font=\scriptsize,red!60!black,align=center] at (0,-3.0) {the client never sees \code{sum.c}};
\end{tikzpicture}
\end{center}

\paragraph{Testing.} It is good practice to write a separate \emph{testing client} for each
module that calls every provided function, so that all tests should pass (in Racket every test
expression should produce \code{\#t}; in C every \code{assert} should succeed).

\section{Designing Modules}

\begin{important}[Goal: high cohesion, low coupling]
\begin{itemize}
  \item \textbf{High cohesion:} all of the interface functions of a module are related and work
    towards a common goal.
  \item \textbf{Low coupling:} there is little interaction between modules. It is impossible to
    remove all coupling, but it should be minimal.
\end{itemize}
\end{important}

\subsection{Information hiding}

A module should hide its implementation details from clients:
\begin{itemize}
  \item \textbf{Security:} we may want to prevent the client from tampering with the data
    used by the module (and breaking its invariants).
  \item \textbf{Flexibility:} if clients cannot depend on implementation details, we are free
    to change the implementation later (e.g.\ for efficiency) without breaking any client.
\end{itemize}

\section{Data Structures and Abstract Data Types}

Consider a module for storing a single account (an ID and a balance). We could implement it
with three different data structures: a two-element list, a structure with two fields, or a
nested structure. If clients only use the interface functions (\code{account-id},
\code{account-balance}, \dots), they cannot tell which one is used.

\begin{definition}
An \emph{abstract data type} (ADT) is a mathematical model for storing and accessing data
through a set of operations. It is described only by the behaviour of its operations, not by
how the data is stored.
\end{definition}

\begin{important}[Data structure vs.\ ADT]
\begin{itemize}
  \item With a \textbf{data structure}, you know how the data is organised and you can access it
    directly in any way you like (e.g.\ a Racket list, a C array).
  \item With an \textbf{ADT}, you do not know how the data is stored; you can only access it
    through the interface functions (operations) the ADT provides.
\end{itemize}
An ADT is \emph{implemented} using some data structure, and the same ADT may have several
implementations with different efficiency.
\end{important}

In Racket it is hard to truly hide a structure from clients. In C we will use \emph{opaque
structures} (Units~06 and~10) to enforce the abstraction.

\section{Collection ADTs}

A \emph{collection ADT} stores an arbitrary number of items and has well-defined operations
for adding, removing and accessing them.

\subsection{Dictionary (revisited)}

The \emph{dictionary} ADT (also called a map, associative array, or symbol table) is a
collection of \emph{key/value pairs} with unique keys. (JSON objects in JavaScript are
dictionaries.) Typical operations:
\begin{itemize}
  \item \code{lookup}: given a key, produce the corresponding value (or ``not found'');
  \item \code{insert}: add a new key/value pair (or replace the value of an existing key);
  \item \code{remove}: delete the pair with a given key.
\end{itemize}
A dictionary can be implemented with an association list, a (sorted) array, or a binary search
tree (as in CS 135 --- we implement it in C in Unit~12).

\subsection{Stack}

A \emph{stack} is a collection where items are ``stacked'' on top of each other. Items are
\emph{pushed} onto the top and \emph{popped} off the top, so the last item in is the first item
out: LIFO. Think of a stack of plates. Uses: browser ``back'' history, undo in text editors, the
call stack (Unit~05), matching brackets.

Operations:
\begin{itemize}
  \item \code{push}: add an item to the top;
  \item \code{pop}: remove the top item;
  \item \code{top}: produce the top item (without removing it);
  \item \code{is\_empty}: determine if the stack is empty.
\end{itemize}

\begin{center}
\begin{tikzpicture}[it/.style={draw,minimum width=1.6cm,minimum height=6mm,fill=blue!6},
  arr/.style={-{Stealth},thick}]
  \node[it] (a) at (0,0) {1};
  \node[it] (b) at (0,0.6) {2};
  \node[it,fill=orange!20] (c) at (0,1.2) {3};
  \draw[thick] (-0.95,1.6) -- (-0.95,-0.4) -- (0.95,-0.4) -- (0.95,1.6);
  \node[font=\scriptsize,anchor=west] at (1.1,1.2) {top};
  \draw[arr] (-2.4,2.4) node[left,font=\scriptsize\ttfamily]{push} to[out=0,in=90] (-0.3,1.6);
  \draw[arr] (0.3,1.6) to[out=90,in=180] (2.4,2.4) node[right,font=\scriptsize\ttfamily]{pop};
\end{tikzpicture}
\end{center}

\paragraph{Stack vs.\ Racket list.} The operations correspond closely to \code{cons},
\code{rest}, \code{first} and \code{empty?}. The significant difference is that only the top
item of a stack is accessible, whereas any element of a list can be accessed (e.g.\ with
\code{list-ref}).

\subsection{Queue}

A \emph{queue} is a collection where items are added at the back and removed from the front:
first in, first out (FIFO), like a line-up at a restaurant. Operations:
\begin{itemize}
  \item \code{add\_back}: add an item to the back;
  \item \code{remove\_front}: remove the front item;
  \item \code{front}: produce the front item;
  \item \code{is\_empty}: determine if the queue is empty.
\end{itemize}

\begin{center}
\begin{tikzpicture}[it/.style={draw,minimum width=9mm,minimum height=7mm,fill=blue!6},
  arr/.style={-{Stealth},thick}]
  \foreach \v [count=\i] in {4,7,2,9} \node[it] (q\i) at (\i*0.9,0) {\v};
  \node[it,fill=orange!20] (q1b) at (0.9,0) {4};
  \draw[arr] (q1.west) -- ++(-1.3,0) node[left,font=\scriptsize\ttfamily]{remove\_front};
  \draw[arr] ($(q4.east)+(1.5,0)$) node[right,font=\scriptsize\ttfamily]{add\_back} -- (q4.east);
  \node[font=\scriptsize,below=1mm of q1] {front};
  \node[font=\scriptsize,below=1mm of q4] {back};
\end{tikzpicture}
\end{center}

\subsection{Sequence}

A \emph{sequence} ADT is useful when you want to retrieve, insert or delete at \emph{any}
position. Operations:
\begin{itemize}
  \item \code{item\_at}: produce the item at a given position;
  \item \code{insert\_at}: insert an item at a given position (later items shift back);
  \item \code{remove\_at}: remove the item at a given position;
  \item \code{length}: the number of items.
\end{itemize}

\begin{aside}
Many languages' ``list'' types (Python's \code{list}, C++'s \code{std::vector}) are really
sequence ADTs. Whether \code{item\_at} is fast depends on the implementation: $\bigO(1)$ with
an array, $\bigO(n)$ with a linked list (Unit~09).
\end{aside}

\section{Summary}
\begin{itemize}
  \item Modules group related functions; the dependency graph is acyclic with a single root.
  \item Interfaces (\code{.h}) contain declarations and documentation; implementations
    (\code{.c}) contain definitions.
  \item C globals have program scope by default; \code{static} makes them module scope.
  \item Aim for high cohesion and low coupling; hide information for security and flexibility.
  \item An ADT is defined by its operations, a data structure by its layout.
  \item Collection ADTs: dictionary, stack (LIFO), queue (FIFO), sequence.
\end{itemize}
